\documentclass[12pt]{article}
\usepackage{amsmath, amssymb, geometry, enumitem}
\geometry{margin=1in}
\setlength{\parskip}{0.6em}
\setlength{\parindent}{0pt}

\begin{document}

\begin{center}
\Large \textbf{ECON 300 -- Labour Economics} \vspace{.5cm}\\
\large \textbf{Quiz 4} \vspace{.5cm}\\
\large \textbf{March 31, 2026}
\end{center}

\begin{enumerate}[label=\textbf{\arabic*.}]

\item In the Oaxaca decomposition, the portion of the wage gap that is attributed to differences in observed characteristics such as schooling and experience is called:

\begin{enumerate}[label=\Alph*)]
    \item the discrimination component
    \item the unexplained component
    \item the explained component
    \item the crowding component
    \item the coefficient bias component
\end{enumerate}

\item Suppose women and men have the same observed productivity characteristics, but women still receive lower wages because firms statistically expect greater labour force interruptions among women. This is best described as:

\begin{enumerate}[label=\Alph*)]
    \item equalizing differences
    \item taste-based discrimination only
    \item statistical discrimination
    \item human capital investment
    \item comparable worth
\end{enumerate}

\item According to the crowding hypothesis, if women are restricted to a narrow set of occupations, then:

\begin{enumerate}[label=\Alph*)]
    \item wages in those occupations will tend to rise because of specialization
    \item wages in those occupations will tend to fall because labour supply is concentrated there
    \item men’s wages in all occupations will necessarily fall
    \item occupational segregation has no effect on wages
    \item wage discrimination disappears in competitive equilibrium
\end{enumerate}
\newpage
\item Which of the following would most likely reduce the unexplained portion of the gender wage gap in an Oaxaca decomposition?

\begin{enumerate}[label=\Alph*)]
    \item adding a variable for years of schooling when schooling differs across groups
    \item removing occupation controls from the regression
    \item increasing the sample size without changing the model
    \item estimating the model without experience
    \item using nominal instead of real wages
\end{enumerate}



\item Which of the following best describes the \textbf{unexplained} component in an Oaxaca decomposition?

\begin{enumerate}[label=\Alph*)]
    \item the portion of the wage gap due to differences in education and experience
    \item the portion of the wage gap due to differences in labour force participation rates
    \item the portion of the wage gap due to differences in observed characteristics only
    \item the portion of the wage gap due to differences in coefficients or returns to characteristics
    \item the portion of the wage gap explained entirely by occupational choice
\end{enumerate}

\end{enumerate}

\section*{ Problem}

Consider the following estimated log annual earnings equations for men and women:

\[
\ln W_m = 1.20 + 0.09S + 0.05EXP
\]

\[
\ln W_f = 0.95 + 0.07S + 0.04EXP
\]

where:

\begin{itemize}
    \item \(W_m\) = annual earnings of men
    \item \(W_f\) = annual earnings of women
    \item \(S\) = years of schooling
    \item \(EXP\) = years of labour market experience
\end{itemize}
\newpage
Suppose the average characteristics of men and women are:

\[
\bar{S}_m = 16, \qquad \bar{EXP}_m = 12
\]

\[
\bar{S}_f = 15, \qquad \bar{EXP}_f = 10
\]

Answer the following:

\begin{enumerate}[label=\textbf{(\alph*)}]

\item Compute the predicted log earnings for men using the male wage equation. \hfill [3 points]

\item Compute the predicted log earnings for women using the female wage equation. \hfill [3 points]

\item Compute the raw gender gap in predicted log earnings:
\[
\ln W_m - \ln W_f
\]
\hfill [2 points]

\item Using the \textbf{male wage structure} as the reference group, compute the portion of the gap explained by differences in characteristics:
\[
(\bar{X}_m - \bar{X}_f)\hat{\beta}_m
\]
\hfill [4 points]

\item Compute the unexplained portion of the gap. \hfill [2 points]

\item Interpret the explained and unexplained components economically. \hfill [1 point]

\end{enumerate}

\end{document}